% !PRINTED 337
\weberpage{337}

and the \gesperrt{Newton approximation method} therefore comes down to this,
that in a first approximation one replaces the curve in the interval
$(\alpha, \beta)$ by the tangent at one of the endpoints, instead of
by its chord, as in the interpolation method.

% SOURCE FIGURE CROP: Bd. I registered scan 363 (printed p. 337), 600 dpi,
% crop x=180,y=900,w=2650,h=950; source-faithful raster, Fig. 21--22.
\weberfigureimage{weber-b1-fig21-22.png}{0.94}

If the second differential quotient vanishes in the interval $(\alpha, \beta)$,
so that the curve has a point of inflection, then the case can
occur that both end tangents lead out of the interval;
in that case the \gesperrt{Newton} method is accordingly not ap-

% SOURCE FIGURE CROP: Bd. I registered scan 363 (printed p. 337), 600 dpi,
% crop x=180,y=2280,w=2650,h=1080; source-faithful raster, Fig. 23--24.
\weberfigureimage{weber-b1-fig23-24.png}{0.94}

plicable (Fig. 21). But it can also happen, when the interval $(\alpha, \beta)$
has already been sufficiently narrowed, that the other case occurs,
in which both tangents lead to interior points of the interval
(Fig. 22). In these cases, however, the regula
falsi will always give a better approximation.

But we shall now assume that $f'(x)$ and $f''(x)$ do not
vanish in the interval $(\alpha, \beta)$. Then the curve does not change
the sense of its curvature, and then certainly \gesperrt{one of the
two end tangents} always has its point of intersection with the
$x$-axis in the interior of the interval. This certainly holds when
the tangent is taken at that endpoint of the interval,

% Weber, Algebra. I.        22
%CONT>>
